\documentclass[11pt,a4paper]{article}

\usepackage[utf8]{inputenc}
\usepackage[T1]{fontenc}
\usepackage{lmodern}
\usepackage[margin=0.76in]{geometry}
\usepackage[table]{xcolor}
\usepackage{array}
\usepackage{tabularx}
\usepackage{booktabs}
\usepackage{enumitem}
\usepackage{titlesec}
\usepackage{fancyhdr}
\usepackage{lastpage}
\usepackage{hyperref}

\definecolor{TextInk}{HTML}{233142}
\definecolor{HeadingBlue}{HTML}{1D3557}
\definecolor{AccentTeal}{HTML}{2A7F92}
\definecolor{AccentGold}{HTML}{C58B4E}
\definecolor{RuleGray}{HTML}{D6E0E5}
\definecolor{TableStripe}{HTML}{F3F8FA}
\definecolor{SoftTint}{HTML}{F7F4EE}
\definecolor{SoftBlue}{HTML}{EEF6F8}
\definecolor{SoftGreen}{HTML}{EEF7F0}
\definecolor{SoftRed}{HTML}{FBF0EF}

\hypersetup{
  colorlinks=true,
  linkcolor=HeadingBlue,
  urlcolor=AccentTeal,
  pdftitle={IB Geography SL Paper 2 Simulation Marking Guide - 2026-05-03},
  pdfauthor={IB Geography SL Preparation}
}

\color{TextInk}
\setlength{\parskip}{0.44em}
\setlength{\parindent}{0pt}
\setlength{\tabcolsep}{5pt}
\renewcommand{\arraystretch}{1.18}
\setlist[itemize]{leftmargin=1.32em,itemsep=3pt,topsep=4pt}
\setlist[enumerate]{leftmargin=1.45em,itemsep=3pt,topsep=4pt}

\newcolumntype{Y}{>{\raggedright\arraybackslash}X}
\newcolumntype{L}[1]{>{\raggedright\arraybackslash}p{#1}}

\titleformat{\section}
  {\normalfont\huge\sffamily\bfseries\color{HeadingBlue}}
  {}{0pt}{}
  [\vspace{0.12em}{\color{AccentGold}\titlerule[1pt]}]

\titleformat{\subsection}
  {\normalfont\Large\sffamily\bfseries\color{AccentTeal}}
  {}{0pt}{}

\titleformat{\subsubsection}
  {\normalfont\large\sffamily\bfseries\color{HeadingBlue}}
  {}{0pt}{}

\titlespacing*{\section}{0pt}{1.05em}{0.45em}
\titlespacing*{\subsection}{0pt}{0.78em}{0.22em}
\titlespacing*{\subsubsection}{0pt}{0.55em}{0.16em}

\pagestyle{fancy}
\setlength{\headheight}{14pt}
\fancyhf{}
\fancyhead[L]{\sffamily\color{AccentTeal}May 3 Marking Guide}
\fancyhead[R]{\sffamily\color{HeadingBlue}IB Geography SL Paper 2}
\fancyfoot[C]{\sffamily\color{HeadingBlue}\thepage\ of \pageref{LastPage}}
\renewcommand{\headrulewidth}{0.4pt}

\newcommand{\term}[1]{\textcolor{HeadingBlue}{\textbf{#1}}}
\newcommand{\exam}[1]{\textcolor{AccentTeal}{\textbf{#1}}}
\newcommand{\score}[1]{\textcolor{AccentTeal}{\textbf{#1}}}
\newcommand{\boxnote}[3]{%
\noindent\colorbox{#1}{%
  \begin{minipage}{0.965\textwidth}
  \sffamily
  \textbf{\textcolor{HeadingBlue}{#2}} #3
  \end{minipage}
}}

\begin{document}

\begin{center}
  {\sffamily\bfseries\color{HeadingBlue}\LARGE May 3 Paper 2 Simulation Marking Guide}\\[0.25em]
  {\sffamily\color{AccentTeal}\large Self-Marking, Error Log, And May 4 Carry-Forward}
\end{center}

\boxnote{SoftRed}{Use after the 75-minute paper only.}{This guide is for
self-marking. It is not an official IB mark scheme. Award marks for accurate
geography, clear source use, valid explanation, named examples, and supported
judgment.}

\section{Strict Marking Rules}

\begin{itemize}
  \item Mark only what was written during the timed simulation. Do not improve
  the answer before scoring.
  \item For \exam{state} and \exam{identify}, give one mark for each correct,
  relevant point.
  \item For \exam{describe}, reward patterns, contrasts, or source evidence.
  Do not require causes unless the answer moves into explanation.
  \item For \exam{explain} and \exam{suggest}, require a plausible chain:
  factor $\rightarrow$ process $\rightarrow$ outcome.
  \item For the 10-mark extended answer, require examples, balance, and a
  judgment that answers the exact command term.
\end{itemize}

\section{Section A Mark Guide}

\subsection{Question 1: Population Distribution -- Changing Population [10]}

\subsubsection{1(a): Oldest Population Structure [1]}

Award \score{1 mark} for \term{Country Beta}. It has the highest median age
(\term{48}) and the highest proportion aged 65 and over (\term{29\%}).

\subsubsection{1(b): Alpha And Beta Contrasts [4]}

Award up to \score{4 marks}: up to \score{2 marks} for each clear contrast
using Source A evidence. Valid contrasts include:

\begin{itemize}
  \item Country Alpha is much more youthful: median age \term{18}, under-15
  share \term{43\%}, and 65+ share \term{3\%}; Country Beta is ageing: median
  age \term{48}, under-15 share \term{12\%}, and 65+ share \term{29\%}.
  \item Fertility is higher in Alpha (\term{4.6}) than in Beta (\term{1.3}).
  \item Alpha is growing rapidly (\term{2.7\%}) while Beta is shrinking
  slightly (\term{-0.3\%}).
  \item Alpha has net out-migration (\term{-1.2\%}) while Beta has net
  in-migration (\term{0.4\%}).
  \item Beta is much more urbanized (\term{92\%}) than Alpha (\term{42\%}).
\end{itemize}

\subsubsection{1(c): Consequences Of Population Structure [4]}

Award up to \score{4 marks}: up to \score{2 marks} for each explained
consequence. Credit either youthful or ageing structure, or one of each.

\begin{itemize}
  \item A youthful population can increase pressure on schools, vaccination,
  maternal healthcare, housing, water, sanitation, and future job creation
  because a high child share raises youth dependency and service demand.
  \item A youthful population can create a future demographic dividend if
  education, health, and jobs allow a large working-age cohort to become
  productive.
  \item An ageing population can increase pension, healthcare, social-care, and
  dependency-ratio pressure because a larger elderly share needs more support.
  \item An ageing population can cause labor shortages or pressure for
  immigration, later retirement, automation, or pronatalist policies.
\end{itemize}

\subsubsection{1(d): Gamma Growth With Replacement-Level Fertility [1]}

Award \score{1 mark} for any valid reason. Good answers include:

\begin{itemize}
  \item net in-migration is positive at \term{0.8\%};
  \item population momentum from a still-sizeable young adult cohort;
  \item life expectancy or lower death rates may keep total population rising.
\end{itemize}

\subsection{Question 2: Global Climate -- Vulnerability And Resilience [10]}

\subsubsection{2(a): Climate Hazards [2]}

Award up to \score{2 marks}, one for each hazard from Source B. Valid answers
include river flooding, storm surge, sea-level rise, saline intrusion, coastal
flooding, heatwaves, drought, or rainfall variability.

\subsubsection{2(b): Delta Coast And High-Income Coastal City [4]}

Award up to \score{4 marks}: up to \score{2 marks} for each difference using
Source B evidence.

\begin{itemize}
  \item The delta coast has a larger population in high-risk zones
  (\term{42\%}) than the high-income coastal city (\term{18\%}).
  \item The delta coast has low to medium adaptive capacity, while the
  high-income coastal city has high adaptive capacity.
  \item The delta coast faces displacement, crop loss, water contamination, and
  disease risk; the high-income city is more associated with expensive
  infrastructure damage, heat stress, and transport disruption.
  \item Both face coastal hazards, but the delta coast is more socially and
  livelihood vulnerable.
\end{itemize}

\subsubsection{2(c): Climate Risk Beyond Physical Hazard [4]}

Award up to \score{4 marks}: up to \score{2 marks} for each explained reason.

\begin{itemize}
  \item Exposure matters: a similar storm or heatwave causes greater impact
  where more people, homes, crops, roads, or assets are located in risky places.
  \item Vulnerability matters: poverty, weak housing, poor health, dependence on
  climate-sensitive farming, or limited insurance can turn a hazard into a
  disaster.
  \item Adaptive capacity matters: warning systems, shelters, drainage,
  healthcare, governance, education, and finance can reduce deaths and losses.
  \item Power and inequality matter: some groups have fewer resources to
  evacuate, adapt, insure, or rebuild.
\end{itemize}

\subsection{Question 3: Global Resource Consumption And Security [10]}

\subsubsection{3(a): Uneven Resource Consumption [1]}

Award \score{1 mark} for any valid indicator from Source C, such as energy use
index, material footprint index, food waste index, or the described
resource-security issues.

\subsubsection{3(b): Resource-Security Contrasts [4]}

Award up to \score{4 marks}: up to \score{2 marks} for each contrast using
Source C evidence.

\begin{itemize}
  \item The high-income consumer economy has much higher energy use
  (\term{90}) than the low-income rural economy (\term{12}).
  \item The high-income consumer economy has a much higher material footprint
  (\term{85}) and food waste index (\term{32}) than the low-income rural
  economy (\term{15} and \term{6}).
  \item The high-income economy's security issue is overconsumption, import
  dependence, and waste; the low-income rural economy's issue is low access,
  biomass reliance, price shocks, and service poverty.
  \item One context is insecure because it consumes too much and relies on
  imports; the other is insecure because people lack reliable access.
\end{itemize}

\subsubsection{3(c): Insecurity From Resource Consumption [4]}

Award up to \score{4 marks}: up to \score{2 marks} for each explained way.

\begin{itemize}
  \item High fossil-fuel consumption increases greenhouse-gas emissions,
  air pollution, and climate risk, which can threaten water, food, health, and
  infrastructure security.
  \item High material consumption can increase mining, deforestation, waste,
  water pollution, habitat loss, and e-waste hazards.
  \item Import dependence can expose countries to price shocks, trade
  disruption, conflict, sanctions, or competition over strategic resources.
  \item Unequal resource access can create social tension, health risks, energy
  poverty, forced migration, or local conflict.
\end{itemize}

\subsubsection{3(d): Stewardship Strategy [1]}

Award \score{1 mark} for a valid stewardship strategy, such as circular
economy, repair and reuse, producer responsibility, renewable energy,
efficiency standards, public transport, water conservation, food-waste
reduction, safer recycling, or fairer access policies.

\section{Section B Mark Guide}

\subsection{Question 4: Visual Stimulus Task [10]}

\subsubsection{4(a): Patterns From Source D [4]}

Award up to \score{4 marks}: up to \score{2 marks} for each described pattern
with evidence.

\begin{itemize}
  \item North Delta has the highest projected population growth
  (\term{38\%}) and also the highest climate-risk index (\term{86}).
  \item Dry Coast has the highest water-stress index (\term{91}), even though
  its climate-risk index (\term{74}) is lower than North Delta and Island Port.
  \item Wealthy Bay has the lowest population growth (\term{8\%}), lowest
  climate-risk index (\term{58}), lowest water-stress index (\term{37}), and
  lowest informal-housing share (\term{3\%}).
  \item Island Port has high climate risk (\term{82}) and water stress
  (\term{72}) but lower projected growth than North Delta or Dry Coast.
  \item Informal housing is highest in North Delta (\term{41\%}) and lowest in
  Wealthy Bay (\term{3\%}).
\end{itemize}

\subsubsection{4(b): Different Human Impacts [4]}

Award up to \score{4 marks}: up to \score{2 marks} for each valid reason.

\begin{itemize}
  \item Adaptive capacity varies: wealthier or better-governed regions may have
  sea walls, shelters, drainage, healthcare, insurance, and emergency planning.
  \item Exposure varies: the number of people, homes, assets, ports, roads, or
  farms in the hazard zone affects total impact.
  \item Housing and infrastructure quality vary: informal housing may be more
  vulnerable to flooding, heat, storm surge, or service failure.
  \item Livelihoods vary: farming, fishing, informal work, tourism, or port
  employment are affected in different ways.
  \item Social inequality, age, health, gender, migration status, and income
  affect the ability to evacuate, adapt, or recover.
\end{itemize}

\subsubsection{4(c): Decision-Maker And Decision [2]}

Award up to \score{2 marks}: \score{1 mark} for a valid decision-maker and
\score{1 mark} for a connected decision.

Valid decision-makers include city planners, national government, emergency
management agencies, water utilities, housing authorities, transport agencies,
NGOs, insurers, or international development banks.

Valid decisions include prioritizing flood protection, informal-housing
upgrading, water conservation, drought planning, evacuation routes, desalination
investment, port resilience, social protection, or infrastructure funding.

\section{Section C Mark Guide}

\subsection{Question 5(a): Climate Resilience Strategies [10]}

Credit any valid examples. Strong answers may use Bangladesh cyclone
preparedness, coastal flood management, early warning systems, heat-action
plans, climate-resilient farming, flood shelters, mangrove restoration,
adaptation finance, or urban green infrastructure.

Good answers should:

\begin{itemize}
  \item define resilience as the ability to prepare for, absorb, recover from,
  and adapt to climate shocks and long-term change;
  \item explain how strategies reduce vulnerability or exposure, not just list
  them;
  \item evaluate success using death reduction, livelihood protection,
  affordability, equity, maintenance, governance, and long-term climate change;
  \item make a judgment about partial or uneven success.
\end{itemize}

Possible judgment: resilience strategies can greatly reduce deaths and losses
when early warning, shelters, planning, and community capacity work together,
but success is uneven where poverty, exposure, weak governance, or rising
climate hazards remain.

\subsection{Question 5(b): Resource Stewardship And Increasing Consumption [10]}

Credit any valid examples. Strong answers may use China waste-import
restrictions and e-waste, Singapore water stewardship, circular economy,
renewable energy, producer responsibility, food-waste reduction, or sustainable
transport.

Good answers should:

\begin{itemize}
  \item define stewardship as responsible management of resources for current
  and future users;
  \item link increasing consumption to environmental and social impacts such as
  emissions, pollution, waste, extraction, water stress, inequality, or import
  dependence;
  \item explain how stewardship strategies reduce demand, improve efficiency,
  reuse materials, protect ecosystems, or improve access;
  \item evaluate limits such as rebound effects, cost, technology dependence,
  weak enforcement, unequal access, or shifting waste to another place;
  \item make a clear \exam{to what extent} judgment.
\end{itemize}

Possible judgment: stewardship can reduce some impacts of rising consumption,
especially where regulation, circular economy, technology, and public behavior
align, but it only partly solves the problem if total consumption keeps rising
or impacts are displaced to poorer regions.

\newpage

\section{Generic 10-Mark Level Guide}

Use this for Question 5. Award the best-fit level, then choose the mark within
the band.

\rowcolors{2}{TableStripe}{white}
\begin{tabularx}{\textwidth}{L{0.14\textwidth}Y}
\toprule
\textbf{Marks} & \textbf{Descriptor} \\
\midrule
0 & No relevant geography, blank answer, or wholly incorrect response. \\
1--2 & Basic general statements. Little or no named example. Mainly
descriptive or asserted. The command term is mostly missed. \\
3--4 & Some relevant knowledge and a partly relevant example. Explanation is
thin, uneven, or list-like. Limited balance or evaluation. \\
5--6 & Clear, mostly accurate answer with named examples and some explanation.
Shows more than one side, but evidence or judgment may be uneven. \\
7--8 & Detailed, well-structured answer with named evidence, geographic
processes, stakeholders or scale, and clear evaluation. A supported conclusion
is present. \\
9--10 & Sustained, precise, and balanced argument. Uses strong case evidence,
clear causal chains, and explicit judgment that answers the exact command term.
Limitations and uneven outcomes are handled convincingly. \\
\bottomrule
\end{tabularx}
\rowcolors{2}{}{}

\section{Command-Term Diagnostics}

\rowcolors{2}{TableStripe}{white}
\begin{tabularx}{\textwidth}{L{0.18\textwidth}Y Y}
\toprule
\textbf{Command} & \textbf{Full-credit behavior} & \textbf{Common penalty} \\
\midrule
State / Identify & Short accurate point, usually source-based & Writing a
long explanation or giving an unsupported guess \\
Describe & Pattern, contrast, trend, or feature with evidence & Explaining
causes before describing what is shown \\
Explain & Cause-and-effect chain using because / therefore logic & Listing
factors without showing how they create an outcome \\
Suggest & Plausible inference based on source evidence and geography & Random
guess with no source link \\
Evaluate & Judge success, limitation, scale, cost, equity, or sustainability &
Only describing a strategy without weighing effectiveness \\
To what extent & Make a degree judgment, not a yes/no answer & No conclusion
or a conclusion that repeats the question \\
\bottomrule
\end{tabularx}
\rowcolors{2}{}{}

\section{May 3 Error Log}

Use this during the 15-minute review block. These rows become the starting
point for the May 4 remediation task.

\begin{tabularx}{\textwidth}{L{0.18\textwidth}Y Y Y}
\toprule
\textbf{Error type} & \textbf{What I wrote or missed} & \textbf{Correct version} & \textbf{May 4 fix} \\
\midrule
Population data & \rule{0pt}{1.35cm} & & \\
Climate vulnerability / resilience & \rule{0pt}{1.35cm} & & \\
Resource security / stewardship & \rule{0pt}{1.35cm} & & \\
Visual stimulus & \rule{0pt}{1.35cm} & & \\
Extended answer evaluation & \rule{0pt}{1.35cm} & & \\
Timing & \rule{0pt}{1.35cm} & & \\
\bottomrule
\end{tabularx}

\section{Carry-Forward Decision}

\begin{tabularx}{\textwidth}{L{0.32\textwidth}Y}
\toprule
\textbf{Prompt} & \textbf{Decision} \\
\midrule
Lowest-scoring Paper 2 unit & \rule{0.62\linewidth}{0.4pt} \\
Lowest-scoring question type & \rule{0.62\linewidth}{0.4pt} \\
One Population item to fix & \rule{0.62\linewidth}{0.4pt} \\
One Climate item to fix & \rule{0.62\linewidth}{0.4pt} \\
One Resources item to fix & \rule{0.62\linewidth}{0.4pt} \\
Exact practice task to redo on May 4 & \rule{0.62\linewidth}{0.4pt} \\
\bottomrule
\end{tabularx}

\vspace{0.5em}

\boxnote{SoftGreen}{May 4 priority rule.}{Redo the exact question type that
lost the most marks. If the problem was visual stimulus, redo Question 4 in
\term{12 minutes}. If the problem was case evidence, rewrite only the Section
C answer with two named examples. If the problem was command terms, rewrite the
first sentence of each answer so it directly matches the command.}

\end{document}
